\section{构建与发布}
\label{sec:build}

仓库里只有三类文件：类、内容、工程脚本。\cref{fig:arch} 已经画过它们的关系，
这里的重点是让「换一台机器」不产生额外心智负担。

\subsection{本地构建}

Windows 上用 \texttt{build.ps1}，macOS 与 Linux 上用 \texttt{Makefile}，两者
都调用 \texttt{latexmk} 完成两遍编译与参考文献。脚本会自动判断语言：英文走 pdfLaTeX，
中文走 XeLaTeX，因此不需要记「这份文件该用哪个引擎」。

\begin{lstlisting}[language=bash,
  caption={常用构建命令。给出 \texttt{-Job} 时只构建其中一份。},
  label={lst:build}]
./build.ps1                 # 两份示例都构建
./build.ps1 -Job main-zh    # 只构建中文示例
./build.ps1 -Clean          # 只清理中间文件
./build.ps1 -Watch          # 改一行编译一次
make en                     # macOS / Linux
\end{lstlisting}

\subsection{持续集成}

CI 做两件事：证明干净机器能编译，以及把 PDF 留档。两个 job 并行跑，英文用
pdfLaTeX、中文用 XeLaTeX 加 Fandol 字体集；任何 push 都会编译，并把成品作为
构建产物（artifact）上传。打 tag 时额外生成一个 Release，附上两份 PDF。

\begin{lstlisting}[language=yaml,
  caption={CI 的关键片段。中文任务通过 \texttt{\textbackslash PassOptionsToClass} 注入 \texttt{fandol} 字体集，避免依赖宿主系统字体。},
  label={lst:ci}]
jobs:
  build:
    strategy:
      matrix:
        include:
          - job: main-en
            engine: pdflatex
          - job: main-zh
            engine: xelatex
    steps:
      - uses: actions/checkout@v4
      - uses: xu-cheng/latex-action@v3
        with:
          root_file: ${{ matrix.job }}.tex
\end{lstlisting}

CI 能保证的只是「编译通过」，它不知道表格有没有溢出、图有没有孤零零占一页。
排版仍然需要人眼过一遍，这一条被写进了发布检查清单。

\subsection{版本}

版本号写进 \verb|\reportsetup{version={...}}|，与 git tag 保持一致。改版时
同时更新 PDF 目视检查、\texttt{CHANGELOG.md} 和 Release 说明——这三件事
任何一件单独做，都会在半年后变成谜。
